% ---------------------------------------------------------------------
\section{Les prix : combien ça coûte, combien ça rapporte}
% ---------------------------------------------------------------------

\begin{astuce}[Exemples recalculés automatiquement]
Tous les montants de ce chapitre sont calculés à la compilation à partir
du fichier \texttt{donnees/parametres.tex}. Changez le prix, le nombre
de pages ou les taux : les tableaux se mettent à jour tout seuls.
\end{astuce}

\subsection{Où va l’argent d’un livre vendu en librairie ?}

Dans la chaîne traditionnelle, la répartition du prix hors taxe d’un
livre est souvent présentée ainsi (ordres de grandeur, très variables
selon les maisons et les genres) :

\begin{center}
\small
\begin{tabularx}{0.82\linewidth}{Lr}
\entete \textbf{Acteur} & \textbf{Part du prix HT} \\
Libraire (remise) & 30 à 40 \% \\
Diffusion et distribution & 15 à 22 \% \\
Fabrication (impression, papier) & 10 à 18 \% \\
Éditeur (frais de structure, marge) & 15 à 20 \% \\
Auteur (droits d’auteur) & 6 à 12 \% \\
\bottomrule
\end{tabularx}
\end{center}

L’IAD ne supprime pas ces maillons : il change surtout la
\textbf{fabrication} (plus chère à l’unité) et supprime le
\textbf{stock} et les \textbf{retours}. Quand vous vous autoéditez,
vous cumulez les parts de l’auteur et de l’éditeur.

\avoir{Wa_bN7TvI58}{\textbf{Pourquoi tes ventes de livres ne te rapportent
rien} : la répartition du prix d’un livre, décortiquée en interview.}

\subsection{Le coût d’impression d’un exemplaire}

En IAD, le coût d’un exemplaire suit presque toujours la même logique :

\begin{center}
\pltitre\large coût = part fixe + (nombre de pages × coût par page)
\end{center}

Pour un roman broché en noir et blanc de 150 à 300 pages, comptez
\textbf{environ 3 à 6~€ HT} l’exemplaire selon la plateforme, le format
et le papier. La couleur peut multiplier ce coût par deux à cinq.

\subsection{Exemple 1 — Vendre sur Amazon avec KDP}

Barème Amazon.fr (broché, encre noire, 110 à 828 pages) : coût
d’impression = \eur{\kdpFixe} + \entier{\exPages} pages × \num{\kdpParPage}~€.
La redevance est de \pct{100*\kdpTaux} du prix HT moins le coût
d’impression, et tombe à \pct{100*\kdpTauxBas} sous \eur{\kdpSeuil}
(règle introduite en 2025).

\begin{center}
\begin{tabularx}{0.82\linewidth}{Lr}
\entete \textbf{Livre de \entier{\exPages} pages vendu \eur{\exPrixTTC} TTC} & \\
Prix hors taxe (TVA \pct{\exTVA}) & \eur{\exPrixHT} \\
Taux de redevance appliqué & \pct{100*\kdpTauxApplique} \\
Part conservée par Amazon & \eur{(1-\kdpTauxApplique)*\exPrixHT} \\
Coût d’impression & \eur{\kdpImpression} \\
\rowcolor{plCreme}\textbf{Redevance par exemplaire} & \textbf{\eur{\kdpRedevance}} \\
\bottomrule
\end{tabularx}
\end{center}
\source{Barème relevé en 2025-2026 ; à vérifier sur le calculateur
officiel KDP, qui fait foi. La redevance est versée par Amazon environ
60 jours après la fin du mois de vente.}

\subsection{Exemple 2 — Vendre en librairie via une plateforme d’IAD}

Quand un libraire commande votre livre via une plateforme distribuée
(BoD, Bookelis, IngramSpark…), il faut lui laisser une remise, et la
plateforme se rémunère sur l’impression et/ou par une commission.

\begin{center}
\begin{tabularx}{0.82\linewidth}{Lr}
\entete \textbf{Livre de \entier{\exPages} pages vendu \eur{\exPrixTTC} TTC} & \\
Prix hors taxe & \eur{\exPrixHT} \\
Remise libraire (\pct{\exRemiseLibraire}) & $-$\,\eur{\libRemise} \\
Commission plateforme / distribution (\pct{\exCommission}) & $-$\,\eur{\libCommission} \\
Impression à la demande & $-$\,\eur{\exCoutIAD} \\
\rowcolor{plCreme}\textbf{Reste pour l’éditeur (et l’auteur)} & \textbf{\eur{\libMarge}} \\
\bottomrule
\end{tabularx}
\end{center}

\begin{attention}[La remise libraire]
En dessous de 30 \% de remise, la plupart des libraires refusent de
commander, car ils perdent de l’argent. Une remise standard en
librairie se situe autour de \textbf{30 à 35 \%} pour une commande à
l’unité, davantage pour une mise en place. Les retours sont rarement
acceptés en IAD : c’est un frein, mais aussi une protection pour vous.
\end{attention}

\subsection{Exemple 3 — IAD ou tirage offset ?}

\begin{center}
\begin{tabularx}{0.82\linewidth}{Lr}
\entete \textbf{Comparaison} & \\
Tirage offset & \entier{\offTirage} ex. pour \eur{\offTotal} \\
Coût unitaire offset & \eur{\offUnitaire} \\
Coût unitaire IAD & \eur{\exCoutIAD} \\
Stockage et logistique (offset) & \eur{\offStockage} / ex. / an \\
\rowcolor{plCreme}\textbf{Ventes à partir desquelles l’offset devient rentable} & \textbf{\entier{\offPointMort} ex.} \\
\bottomrule
\end{tabularx}
\end{center}

Autrement dit : si vous n’êtes pas raisonnablement sûr de vendre plus
d’environ \entier{\offPointMort} exemplaires, l’IAD est moins risqué —
sans compter le stockage, les retours et le pilon évités. Au-delà,
un tirage (offset ou numérique) devient plus intéressant, et l’IAD peut
prendre le relais une fois le stock écoulé.

\subsection{Les autres coûts à prévoir}

\begin{center}
\small
\begin{tabularx}{\linewidth}{P{4.5cm}LP{3.2cm}}
\entete \textbf{Poste} & \textbf{Commentaire} & \textbf{Ordre de grandeur} \\
Frais de mise en ligne & gratuits chez KDP et IngramSpark ; offres gratuites
  ou forfaits payants chez BoD, Bookelis et d’autres & 0 à quelques centaines d’€ \\
Frais annuels & certaines offres facturent un abonnement par titre & 0 à quelques dizaines d’€ / an \\
Épreuves et exemplaires auteur & au coût d’impression + port & 5 à 15~€ l’exemplaire livré \\
Correction, mise en page & à ne pas négliger : c’est ce qui fait un livre pro & 300 à 2\,000~€ \\
Couverture & illustrateur ou graphiste & 150 à 1\,000~€ \\
ISBN, dépôt légal & gratuits en France & 0~€ \\
\bottomrule
\end{tabularx}
\end{center}
\source{Fourchettes indicatives observées sur le marché français ; à
ajuster selon vos devis.}
